\chapter{Implementierung und Integration}
\label{chap:implementierung}
In diesem Kapitel wird die Umsetzung der verschiedenen Komponenten des Projekts beschrieben. Es wird erläutert, wie die Hardware und Software implementiert wurden, um die gewünschten Funktionen zu realisieren.
\section{Hardware}

\section{Software}

    \subsection{Websocket}
    % Um die Kommunikation zwischen dem Backend und der Weboberfläche zu ermöglichen, wird ein Websocket verwendet. Dieser wird auf dem Raspberry Pi implementiert und ermöglicht der Weboberfläche, Sensordaten und Kamerabilder zu empfangen und Steuerbefehle an das Backend zu senden. Darüber soll das Fahrzeug über die Weboberfläche gesteuert werden können. Durch die Verwendung von Websockets kann eine Echtzeit-Kommunikation zwischen der Weboberfläche und dem Backend gewährleistet werden, was für die Steuerung und Überwachung des Fahrzeugs von entscheidender Bedeutung ist.
    Es wird die Python-Bibliothek \textbf{PySide6} und die darin enthaltene \textbf{QtWebSockets} Klasse verwendet, um die Websocket-Verbindung zu implementieren. Auf der Weboberfläche wird JavaScript verwendet, um die Verbindung zum Websocket herzustellen und die Daten zum Senden und Empfangen zu verarbeiten bzw. vorzubereiten.
    % noch schreiben, dass Websocket in lokalem Netzwerk läuft und feste IP-Adresse hat, damit die Weboberfläche immer weiß, wie sie das Backend erreichen kann?
    %noch erwähnen, dass weboberfläche auf laptop läuft?

    \subsection{Videoübertragung} % wie aufgenommen? wie gesendet? etc.
    Das Aufnehmen der Bilder erfolgt mittels der Raspberry Pi Kamera, die über die CSI-Schnittstelle mit dem Raspberry Pi verbunden ist. Jedoch ist das Aufnehen und die Übertragung dieser Bilder kein triviales Problem. Die Bilder müssen in Echzeit aufgenommen, komprimiert und über das Netzwerk an den Client gesendet werden. Wie zuvor im Abschnitt (encoding vergleich ref) festgelegt, soll als Übertragungsprotokoll WebRTC und als Codec H.264 verwendet werden, da es speziell für die Übertragung in Echtzeit entwickelt wurde und eine geringe Latenz bietet. 

    Eine eigene Implementierung des WebRTC Protkolls wäre sehr aufwendig und würde nicht zum eigentlichen Ziel der Arbeit beitragen. Auf Grund dessen wird die Software \textit{MediaMTX} verwendet, welche eine einfache Schnittstelle für die Implementierung von WebRTC bietet. MediaMTX ermöglicht es, die aufgenommenen Bilder der Kamera zu komprimieren und über das Netzwerk an den Client zu senden, ohne dass eine eigene Implementierung des WebRTC Protokolls erforderlich ist. Ein weiterer Vorteil von \textit{MediaMTX} ist, dass es im Falle er PI Kamera direkt die Bilder von der Kamer aufnehmen und komprimieren kann, ohne dass die Bilder zuerst an die CPU gesendet werden müssen, was die Latenz weiter reduziert. Somit profitiert die Videoübertragung von den Hardwareencodern des Raspberry Pi, was zu einer deutlich geringeren Latenz führt.

    Damit die Kamera von MediaMTX erkannt und verwendet werden kann, muss sie in der Konfigurationsdatei von MediaMTX entsprechend konfiguriert werden. In Listing~\ref{lst:mediamtx_config} ist die Konfiguration der Kamera in der MediaMTX Konfigurationsdatei zu sehen. Hier werden die Quelle der Bilder, die Auflösung, die Bildrate und die Ausrichtung der Bilder festgelegt. Durch diese Konfiguration kann MediaMTX die Bilder der Kamera aufnehmen, komprimieren und über das Netzwerk an den Client senden. Auf Grund dessen, dass die Kamera verkehrt herum montiert ist, müssen die Bilder vertikal und horizontal gespiegelt werden, damit sie auf der Weboberfläche korrekt angezeigt werden.

    \begin{lstlisting}[caption={Konfiguration der Kamera in der MediaMTX Konfigurationsdatei}, label={lst:mediamtx_config}]
    cam:
        source: rpiCamera
        rpiCameraWidth: 1280
        rpiCameraHeight: 720
        rpiCameraFPS: 30
        rpiCameraVFlip: yes
        rpiCameraHFlip: yes
    \end{lstlisting}

    Die Weboberfläche kann sich mit dem WebRTC Stream von MediaMTX verbinden und die Bilder in einem Video Element anzeigen. 

    \subsection{Notbremsung} % wir könnten den Scheiß richtig kompliziert machen, wenn wir berechnen wie schnell sich der Abstand zum Hindernis reduziert, wie schnell das Fahrzeug ist und wie weit das Hindernis entfernt ist. Weil wenn man hinter einem fahrenden Auto fährt und man nur Geschwindigkeit und Abstand misst, wird die Notbremsung vielleicht ungewollt ausgelöst. Oder wir sagen wir beachten nur 
    In diesem Projekt werden Ultraschallsensoren als Frühwarnsystem zur Hinderniserkennung eingesetzt. Die Sensoren werden so positioniert, dass sie den Bereich vor dem Fahrzeug überwachen können. Wenn ein Hindernis erkannt wird und die Gefahr einer Kollision besteht, kann automatisch eine Notbremsung eingeleitet werden, um diese zu verhindern. \\
    Ultraschallsensoren sind für diesen Anwendungsfall besonders geeignet, da sie zuverlässig Entfernungen messen können, selbst wenn die Sichtverhältnisse schlecht sind, beispielsweise bei Dunkelheit, Staub oder Nebel. Zudem sind sie kostengünstig und einfach zu integrieren, was sie zu einer idealen Wahl für die Hinderniserkennung in diesem Projekt macht. \\

    Eine Notbremsung wird durchgeführt, wenn ein Hindernis eine potenzielle Kollisionsgefahr darstellt. Die gemessene Entfernung zu diesem Hindernis wird unter Berücksichtigung der aktuellen Geschwindigkeit des Fahrzeugs ausgewertet. Dadurch kann ermittelt werden, ob eine Kollision wahrscheinlich ist. \\
    Die Auswertung der gemessenen Entfernung erfolgt durch die Berechnung des Bremswegs, der anhand der aktuellen Geschwindigkeit berechnet wird. Hierzu wird die Formel für eine Notbremsung herangezogen. Die Formel für die Berechnung des Bremswegs bei einer normalen Bremsung kann in diesem Fall vernachlässigt werden, da es sich um die Implementierung einer Notbremsung handelt. Für eine Notbremsung wird folgende Formel verwendet: \\
    \begin{equation}
        s = \frac{v}{10} \cdot \frac{v}{10} \cdot 0.5
    \end{equation} 
    wobei \( s \) dem Bremsweg in Metern und \( v \) der Geschwindigkeit in km/h entspricht. \\
    Dabei wird der Bremsweg um 50 \% reduziert, um die schnellere Reaktion und stärkere Betätigung des Bremspedals bei einer Notbremsung zu berücksichtigen \cite{BremswegBerechnen}. 
    
    Für gewöhnlich wird für den tatsächlichen Anhalteweg eines Fahrzeugs zusätzlich zum Bremsweg auch die Reaktionszeit berücksichtigt, die der Fahrer benötigt, um auf die Gefahr zu reagieren und die Bremsung einzuleiten. Da in diesem Fall jedoch eine Notbremsung durch die Software automatisch ausgelöst wird, fällt diese Reaktionszeit deutlich geringer aus. Ein weiterer Aspekt der berücksichtigt werden muss, ist die Tatsache, dass das Fahrzeug aufgrund seines geringen Gewichts einen deutlich kürzeren Bremsweg hat als ein echtes Fahrzeug. Daher wird die Differenz zwischen dem Bremsweg eines echten Fahrzeugs und diesen Fahrzeugs als Puffer für die Reaktionszeit und die Bremsverzögerung genutzt, um sicherzustellen, dass die Notbremsung nicht zu spät ausgelöst wird.

    %!!!! Fahrzeug hat viel kürzeren Bremsweg, weil kein Gewicht
    % aber wir nehmen einfach die normale Formel für Notbremsung und sagen, dass die Differenz zwischen Bremsung eines echten Fahrzeug und unseres Fahrzeus als Puffer für Reaktionszeit und Bremsverzögerung etc. dient, damit die Notbremsung nicht zu spät ausgelöst wird.

    \subsection{Geschwindigkeitsmessung} % oben hall sensor kapitel raus?
    Um die Geschwindigkeit des Fahrzeugs zu messen, wird ein Hall-Sensor verwendet, der die Drehung eines Rades erfasst. Dabei werden Magneten im Rad verbaut, welche vom Hall-Sensor erkannt werden. Sobald ein Magnet in der Nähe des Sensors ist, wird ein Signal generiert. Durch die Messung der Zeit, die für eine vollständige Umdrehung des Rades benötigt wird und dem bekannten Radumfang, kann die Geschwindigkeit des Fahrzeugs berechnet werden. \\

    \subsection{Schilderkennung}
    \subsection{Fernsteuerung}
    \subsection{Follow the line}
    \subsection{etc.}
